\chapter{热力学基础}
在\xref{chap:气体动理论}中我们主要讨论了理想气体在平衡状态下的基本性质和统计规律，而本章将讨论理想气体在受到外界影响时其状态变化所遵循的规律。实验表明，在热力学的概念适用的范围内，一切与热运动有关的现象都遵循热力学基本定律，这包含热力学第一定律和第二定律
\begin{enumerate}
    \item 热力学第一定律是关于能量的规律，是获得物质各种宏观性质的依据。
    \item 热力学第二定律是关于熵的原理，是判断宏观方程进行方向和限度的规律。
\end{enumerate}
本章将重点介绍热力学定律在热力学过程中的应用，并简要分析熵原理的统计意义。

\input{Chapter02A.tex}
\input{Chapter02B.tex}
\input{Chapter02C.tex}
\input{Chapter02D.tex}
\input{Chapter02E.tex}